% Lecture notes: Week 6, Lecture 1 -- Shear and Moment Diagrams
% To hide this lecture from the website, add a line containing only:  % draft
% To set the date shown on the website, add a line like:  % posted: 2026-09-02
\documentclass[11pt]{article}
\usepackage{mech2030notes}
\begin{document}

% ##################################################################
%  WEEK 6, LECTURE 1
% ##################################################################
\lecture{6}{1}{Shear and Moment Diagrams}

\noindent
Shear and moment diagrams are \emph{plots of the internal shear $V(x)$ and
internal moment $M(x)$ along a beam}.
\begin{itemize}[label=$\rightarrow$]
  \item The \emph{maximum bending moment} $M_{\max}$ corresponds to the
        \emph{maximum bending stress} (more on this later).
  \begin{itemize}[label=\then]
    \item These diagrams tell us \emph{where} $M_{\max}$ occurs.
  \end{itemize}
\end{itemize}

% ==================================================================
\section*{Example: Squirrel Problem}
% ==================================================================

A squirrel of weight $\SI{5}{N}$ sits at the tip of a branch, modeled as a
cantilever beam of length $\SI{0.4}{m}$ fixed at $A$.

\begin{center}
\begin{tikzpicture}
  \wall{0.9}
  \draw[beam] (0,-0.15) rectangle (6,0.15);
  \draw[force] (6,1.4) node[above] {$\SI{5}{N}$} -- (6,0.17);
  \node[left] at (-0.4,0) {$A$};
  \draw[dim] (0,-0.7) -- (6,-0.7) node[midway, fill=white] {$\SI{0.4}{m}$};
  \draw[axis] (0,-1.35) -- (1.1,-1.35) node[right] {$x$};
  \draw[thin] (0,-1.2) -- (0,-1.5);
\end{tikzpicture}
\end{center}

\subsection*{Steps to draw $V(x)$ and $M(x)$}

% ------------------------------------------------------------------
\subsubsection*{1) Solve the global FBD for the reactions}

\begin{center}
\begin{tikzpicture}
  \draw[thick] (0,0) -- (6,0);
  \node[pin] at (0,0) {};
  \node[pin] at (6,0) {};
  % reactions
  \draw[force] (-1.4,0) node[left] {$A_x$} -- (-0.08,0);
  \draw[force] (0,-1.3) node[below] {$A_y$} -- (0,-0.08);
  \draw[force] ([shift={(20:0.6)}]0,0) arc (20:160:0.6);
  \node at (0,0.95) {$M_A$};
  % applied load
  \draw[force] (6,1.3) node[above] {$\SI{5}{N}$} -- (6,0.08);
  \draw[dim] (0,-2.2) -- (6,-2.2) node[midway, fill=white] {$\SI{0.4}{m}$};
  \draw[thin, gray] (6,-0.15) -- (6,-2.35);
\end{tikzpicture}
\end{center}

\begin{align*}
  \textstyle\sum F_x = 0 \quad &\Rightarrow \quad \boxed{A_x = 0} \\
  \textstyle\sum F_y = 0 \quad &\Rightarrow \quad \boxed{A_y = \SI{5}{N}} \\
  \textstyle\sum M^{(A)} = 0 \quad &\Rightarrow \quad
      M_A - (\SI{0.4}{m})(\SI{5}{N}) = 0
      \quad\Rightarrow\quad \boxed{M_A = \SI{2.0}{N.m}}
\end{align*}

% ------------------------------------------------------------------
\subsubsection*{2) Use the key equations for each load type}

\paragraph{2a) Distributed loads $p(x)$:} (from HW\,2)
\[
  \boxed{\frac{dV}{dx} = p(x)}
  \qquad\qquad
  \boxed{\frac{dM}{dx} = V(x)}
\]

\paragraph{2b) Point loads and point moments:}
Consider a point load $P_0$ and a point moment $M_0$ acting on a beam, and
take a thin element of width $\Delta x$ around each.

\begin{center}
\begin{tikzpicture}
  \wall{0.9}
  \draw[beam] (0,-0.15) rectangle (6,0.15);
  % point load element
  \draw[element] (2.2,-0.45) rectangle (2.8,0.45);
  \draw[force] (2.5,0.17) -- (2.5,1.3) node[above] {$P_0$};
  % point moment element
  \draw[element] (4.7,-0.45) rectangle (5.3,0.45);
  \node[pin] at (5,0) {};
  \draw[force] ([shift={(-40:0.32)}]5,0) arc (-40:200:0.32);
  \node at (5,0.75) {$M_0$};
\end{tikzpicture}
\end{center}

% --- element free-body diagrams ---
\noindent
\begin{minipage}[b]{0.48\textwidth}
\centering
\textbf{Point load}\\[0.6em]
\begin{tikzpicture}
  \draw[thick] (0,0) rectangle (1.6,2);
  \draw[force] (-0.3,0.2) -- (-0.3,1.8);
  \node[left] at (-0.35,1) {$V(x)$};
  \draw[force] (1.9,1.8) -- (1.9,0.2);
  \node[right] at (1.95,1) {$V(x+\Delta x)$};
  \draw[force] (0.8,2.05) -- (0.8,3.0) node[above] {$P_0$};
  \draw[dim] (0,-0.4) -- (1.6,-0.4) node[midway, below=2pt] {$\Delta x$};
\end{tikzpicture}
\end{minipage}%
\hfill
\begin{minipage}[b]{0.48\textwidth}
\centering
\textbf{Point moment}\\[0.6em]
\begin{tikzpicture}
  \draw[thick] (0,0) rectangle (1.6,2);
  \node[pin] at (0.8,1) {};
  % M(x): clockwise on left face
  \draw[force] ([shift={(235:0.9)}]0,1) arc (235:125:0.9);
  \node[left] at (-0.95,1) {$M(x)$};
  % M(x+dx): counterclockwise on right face
  \draw[force] ([shift={(-55:0.9)}]1.6,1) arc (-55:55:0.9);
  \node[right] at (2.55,1) {$M(x+\Delta x)$};
  % M0: counterclockwise at center
  \draw[force] ([shift={(-40:0.35)}]0.8,1) arc (-40:200:0.35);
  \node at (0.8,1.65) {$M_0$};
  \draw[dim] (0,-0.4) -- (1.6,-0.4) node[midway, below=2pt] {$\Delta x$};
\end{tikzpicture}
\end{minipage}

% --- equilibrium ---
\noindent
\begin{minipage}[t]{0.48\textwidth}
\begin{gather*}
  \textstyle\sum F_y = 0 \quad (+\uparrow) \\
  V(x) + P_0 - V(x+\Delta x) = 0 \\[2pt]
  \boxed{V(x+\Delta x) = V(x) + P_0}
\end{gather*}
\end{minipage}%
\hfill
\begin{minipage}[t]{0.48\textwidth}
\begin{gather*}
  \textstyle\sum M = 0 \quad (+\circlearrowleft) \\
  -M(x) + M_0 + M(x+\Delta x) = 0 \\[2pt]
  \boxed{M(x+\Delta x) = M(x) - M_0}
\end{gather*}
\end{minipage}

\vspace{0.8em}
% --- resulting jumps ---
\noindent
\begin{minipage}[t]{0.48\textwidth}
\centering
\begin{tikzpicture}
  \draw[axis] (0,0) -- (4,0) node[right] {$x$};
  \draw[axis] (0,0) -- (0,2.8) node[above] {$V(x)$};
  \draw[plotline] (0,1.0) -- (2,1.0);
  \draw[plotline] (2,2.1) -- (3.6,2.1);
  \draw[force] (2,1.02) -- (2,2.08);
  \node[right] at (2.05,1.55) {$P_0$};
  \node[below] at (1,1.0) {$V(x)$};
  \node[above] at (2.8,2.1) {$V(x+\Delta x)$};
\end{tikzpicture}\\[0.3em]
{\small Shear jumps \emph{up} by $P_0$.}
\end{minipage}%
\hfill
\begin{minipage}[t]{0.48\textwidth}
\centering
\begin{tikzpicture}
  \draw[axis] (0,0) -- (4,0) node[right] {$x$};
  \draw[axis] (0,0) -- (0,2.8) node[above] {$M(x)$};
  \draw[plotline] (0,2.1) -- (2,2.1);
  \draw[plotline] (2,1.0) -- (3.6,1.0);
  \draw[force] (2,2.08) -- (2,1.02);
  \node[right] at (2.05,1.55) {$M_0$};
  \node[above] at (1,2.1) {$M(x)$};
  \node[below] at (2.8,1.0) {$M(x+\Delta x)$};
\end{tikzpicture}\\[0.3em]
{\small Moment jumps \emph{down} by a counterclockwise $M_0$.}
\end{minipage}

% ------------------------------------------------------------------
\subsubsection*{3) Shear diagram}

\noindent
\begin{minipage}[c]{0.48\textwidth}
\centering
\textbf{FBD}\\[0.4em]
\begin{tikzpicture}
  \draw[thick] (0,0) -- (4,0);
  \node[pin] at (0,0) {};
  \node[pin] at (4,0) {};
  \draw[force] ([shift={(20:0.55)}]0,0) arc (20:160:0.55);
  \node[above] at (0,0.6) {$\SI{2}{N.m}$};
  \draw[force] (0,-1.2) node[below] {$\SI{5}{N}$} -- (0,-0.08);
  \draw[force] (4,1.2) node[above] {$\SI{5}{N}$} -- (4,0.08);
  \draw[dim] (0,-2.1) -- (4,-2.1) node[midway, fill=white] {$\SI{0.4}{m}$};
  \draw[thin, gray] (4,-0.1) -- (4,-2.25);
\end{tikzpicture}
\end{minipage}%
\hfill
\begin{minipage}[c]{0.48\textwidth}
\centering
\begin{tikzpicture}
  % scale: 0.4 m -> 4 cm, 5 N -> 2.5 cm
  \fill[hatch] (0,0) rectangle (4,2.5);
  \draw[axis] (0,0) -- (5,0) node[right] {$x$};
  \draw[axis] (0,0) -- (0,3.3) node[above] {$V(x)$};
  \draw[plotline] (0,0) -- (0,2.5) -- (4,2.5) -- (4,0);
  \node[pin] at (0,0) {}; \node[pin] at (0,2.5) {};
  \node[pin] at (4,2.5) {}; \node[pin] at (4,0) {};
  \node[left] at (0,2.5) {$\SI{5}{N}$};
  \node[below] at (4,0) {$\SI{0.4}{m}$};
  \draw[force, fill=white] (0.35,0.15) -- (0.35,2.35);
  \node[fill=white, inner sep=1.5pt, right] at (0.4,1.25) {$+A_y$};
  \draw[force] (3.65,2.35) -- (3.65,0.15);
  \node[fill=white, inner sep=1.5pt, left] at (3.6,1.25) {$-\SI{5}{N}$};
\end{tikzpicture}
\end{minipage}

\begin{itemize}
  \item Point load of $+\SI{5}{N}$ at $x = 0$.
  \item $p(x) = 0$ on $0 < x < \SI{0.4}{m}$, so $\dfrac{dV}{dx} = 0$
        (constant shear).
  \item Point load of $-\SI{5}{N}$ at $x = \SI{0.4}{m}$.
\end{itemize}

% ------------------------------------------------------------------
\subsubsection*{4) Moment diagram}

\begin{itemize}
  \item Point moment of $+\SI{2}{N.m}$ at $x = 0$.
  \item $V(x) = \SI{5}{N}$, so
        $\dfrac{dM}{dx} = V = +\SI{5}{N.m/m}$.
\end{itemize}

\begin{center}
\begin{tikzpicture}
  % scale: 0.4 m -> 4 cm, 1 N.m -> 1 cm
  \fill[hatch] (0,0) -- (0,-2) -- (4,0) -- cycle;
  \draw[axis] (0,0) -- (5,0) node[right] {$x$};
  \draw[axis] (0,-2.4) -- (0,1) node[above] {$M(x)$};
  \draw[plotline] (0,0) -- (0,-2) -- (4,0);
  \node[pin] at (0,0) {}; \node[pin] at (0,-2) {}; \node[pin] at (4,0) {};
  \node[above] at (4,0) {$\SI{0.4}{m}$};
  \node[left] at (0,-2) {$-\SI{2}{N.m}$};
  % jump due to point moment
  \draw[force] (-0.4,-0.05) -- (-0.4,-1.7);
  \node[left] at (-0.45,-0.9) {$-M_A$};
  % slope triangle
  \draw[thin] (2,-1) -- (3,-1) -- (3,-0.5);
  \node[below right] at (2.5,-1) {$\SI{5}{N.m/m}$};
\end{tikzpicture}
\end{center}

\begin{itemize}[label=\then]
  \item Here, the largest moment (in magnitude) occurs at $x = 0$.
  \begin{itemize}[label=\then]
    \item The largest stress is at $x = 0$, so the branch will break there.
  \end{itemize}
\end{itemize}

\end{document}
